\documentclass[11pt]{article}
\usepackage{lmodern}
\usepackage[T1]{fontenc}
\usepackage[letterpaper,margin=0.8in]{geometry}
\usepackage{xcolor}
\usepackage{booktabs}
\usepackage{tabularx}
\usepackage{longtable}
\usepackage{enumitem}
\usepackage{array}
\usepackage[hidelinks]{hyperref}

\setlist{itemsep=1pt,topsep=3pt}
\newcommand{\br}[1]{\texttt{#1}}
\newcommand{\h}[1]{\texttt{#1}}
\newcolumntype{L}{>{\raggedright\arraybackslash}X}

\title{PACMAN Legacy Firmware (1.X)\\[2pt]
  \large Branches, build steps and bitwise reproduction}
\author{Maintained by the PACMAN agents (augie)}
\date{2026-10-31}

\begin{document}
\maketitle
\sloppy

\section{Purpose and scope}

PACMAN firmware comes in three orthogonal families: \emph{legacy} (1.X, everything before consistent versioning),
\emph{hijinks} (2.X, the 2024 firmware for the 1v5 test stands, described in a companion note) and \emph{zipline}
(3.X, the current firmware in the \h{pacman-firmware} and \h{pacman-software} repositories).  This note describes
how the legacy firmware is kept in the repository \h{mulhearn/pacman}: one branch per hardware line, each holding
as much of that line's firmware as possible, every state that was released or deployed reproducible from the
repository alone, and the proof of each reproduction.

A bitstream is \emph{reproduced bitwise} when the SHA-256 of its configuration body (the bitstream without its
header, which carries the build date and paths) equals that of the reference bitstream.  The bodies are quoted here
by their first eight hex digits.

\section{Layout}

\begin{itemize}
\item One branch per hardware line (\br{hw1v3}, \br{hw1v4}, \br{hw1v5}), firmware only, as linear as the history
  allows.
\item Historic commits keep their original hashes.  New commits go only at a branch's tip:
  \begin{itemize}
  \item \emph{snapshots}: firmware that was built outside the repository, restored from its Vivado project archive;
  \item \emph{build steps}: what makes a state build from the repository alone (fixes without behavior change,
    staging scripts); they stay on the branch;
  \item \emph{build options}: variants of a build (rot8).
  \end{itemize}
\item Superseded side lines, which change nothing on the main lines, are recorded on two \emph{excursions}
  branches, firmware and software.
\item Branches under review are named \br{working-on-}\emph{name} until they are frozen (renamed without the
  prefix and tagged); they may be rewritten only until then.
\item A checker (\h{check-legacy-hw1vN.sh}, run on fresh clones) enforces the layout: each spine ends at its
  historic tip with the original hashes and is followed by exactly the listed new commits, the build steps change
  only build scripts, the excursions never change the tree, and commit messages cite only commits on the spines.
\end{itemize}

\begin{table}[h]
\centering\small
\begin{tabularx}{\textwidth}{@{}lllL@{}}
\toprule
Branch & Tip & Historic tip & New at the tip \\
\midrule
\br{working-on-legacy-firmware-hw1v3} & 9fbc1dc & 501180f & 893 snapshot, build step, 883 snapshot \\
\br{working-on-legacy-firmware-hw1v4} & d08d199 & 21019ee & 894 (two commits), build step \\
\br{working-on-legacy-firmware-hw1v5} & 7062e17 & 2425e67 & SYNC FIX build step, build option rot8 \\
\br{working-on-legacy-firmware-excursions} & d5fd1c7 & (2425e67) & three record merges \\
\br{working-on-legacy-software-excursions} & 4a27ae1 & (2425e67) & three record merges \\
\bottomrule
\end{tabularx}
\caption{The legacy branches (2026-10-31).}
\end{table}

The first commit of all of them is 0b9bee4 (2019-11-16, the Trenz TE0720 skeleton).  Zipline shares only that
commit with legacy; \h{pacman-firmware} and \h{pacman-software} share no history with this repository.

\section{The hw1v3 line}

\paragraph{Spine.}  The trunk (v1.2, v1.2b, three 2020 merges), then the v1.3 releases:  v1.3-0.0 (a10a44a,
2020-12-23, the first v1.3 board), v1.3-0.1 (e26918e, external sync and trigger inputs), v1.3-0.2 (1dfc351, no
firmware change), v1.3-0.3 (2099d0d, POSI/PISO invertible), and v1.3-1.0 (501180f, 2022-05-16, Peter Madigan's
adjustable MCLK).  Every tag builds as committed with Vivado 2018.3 (the first run fails at IP generation, the
rerun builds).

\paragraph{New commits.}
\begin{itemize}
\item \textbf{893} (a5f8e91): Armin Karcher's V3 firmware (2024-02), the version that runs on the 1v3 boards
  (version register \h{0x37D}), from its project archive \h{pacman-fw893.xpr.zip}: the archive's project Tcl as
  \h{tcl/pacman-fw15.tcl} (Vivado 2021.1) and the sources it needs that differ from 501180f
  (\h{clk\_ddr\_en.vhd}, \h{uart\_en.vhd} new; \h{larpix\_mclk\_sel.vhd}, \h{larpix\_reset\_gen.vhd} changed; the
  FIFO \h{.xci} as upgraded to 2021.1).
\item \textbf{Build step} (1f66dcd): \h{tcl/build\_fw15.sh} and \h{tcl/build\_fw15.tcl}.
\item \textbf{883} (9fbc1dc): 893 with edge-sensitive SYNC/TRIG inputs (Armin, 2024-05, for the 2x2 trigger
  module): \h{larpix\_reset\_gen.vhd} and that archive's project Tcl.
\end{itemize}

\paragraph{Build.}  From the top of a checkout, with the Vivado 2021.1 settings sourced:
\begin{verbatim}
bash tcl/build_fw15.sh        # -> products/zsys_wrapper.bit
\end{verbatim}
The project Tcl is Armin's export, unchanged: it names its files in the layout of his Vivado project and has no
top-level wrapper.  The script stages the files it requires (its own \h{checkRequiredFiles} list) from
\h{srcs/}, \h{constraints/} and \h{srcs/ip} into that layout under \h{build/fw15}, runs the export, adds the
wrapper as \h{tcl/build\_xpr.tcl} does, and builds through \h{write\_bitstream}.

\section{The hw1v4 line}

\paragraph{Spine.}  From v1.3-0.3: the rev4 commits (new constraints with one GLOBAL\_SYNC, I2C virtual
registers, eMMC), v1.4-0.0 (4532a23, 2022-04-12), the merge of v1.3-1.0 (b37e07f), v1.4-0.2 (0681491,
2022-10-04, the last larpix release), and Armin's module2 upload (af8c85b, 21019ee, 2023-03-01: the resethack 1.8
project Tcl \h{tcl/pacman\_resethack1.8.tcl} and \h{pacman\_r4.xdc}).

\paragraph{New commits.}
\begin{itemize}
\item \textbf{894}, resethack 1.8 (a693c5e, c98f6b0): the sources the module2 Tcl needs that were never
  committed (\h{uart\_en.vhd}, \h{tristate\_out.vhd}, \h{clk\_ddr\_en.vhd}, \h{larpix\_reset\_gen.vhd} with the
  per-tile reset mask, \h{zsys\_wrapper.vhd}), from the project archive \h{pacman\_resethack1.85.xpr.zip}, and the
  FIFO \h{.xci} as upgraded to 2021.1.
\item \textbf{Build step} (d08d199): \h{tcl/build\_resethack1.8.sh}, \h{tcl/recreate\_resethack1.8.tcl},
  \h{tcl/build\_resethack1.8.tcl}.
\end{itemize}

\paragraph{Build.}  With the Vivado 2021.1 settings sourced:
\begin{verbatim}
bash tcl/build_resethack1.8.sh     # -> products/zsys_wrapper.bit
\end{verbatim}
The recreate script stages every file Armin's Tcl requires under \h{build/resethack1.8}, links the Windows path
by which it imports the block design (relative on Linux) to the staged tree, and sources a copy with one edit:
Vivado 2021.1 places \h{pacman\_r4.xdc} in \h{constrs\_1/imports/constraints/}, so the target constraints path is
changed to match.  The build launches the implementation a second time when the first launch stops at the 2021.1
IP generation.

\section{The hw1v5 line}

\paragraph{Spine.}  From v1.4-0.0: the 2024 work (c2872b5, new constraints, IP upgrade, Vivado 2020.2; 36e32af,
SYNC FIX: RESETN replaced by a per-tile SYNC; 5588b58, constraints cleanup; d8790e0, rot8 constraints as a
separate file), ending at 2425e67 (2024-08-11).  2425e67 does not build as committed.

\paragraph{New commits.}
\begin{itemize}
\item \textbf{Build step} (3155d4a, ``restore build from repo''): four fixes without behavior change:
  \h{build\_xpr} writes the hardware description with \h{write\_hwdef} (Vivado 2020.2 writes no \h{.sysdef});
  the data\_tx subset converter's options are set after \h{validate\_bd\_design} (AR 58038);
  \h{larpix\_mclk\_sel} has integer generics, so the block design regenerates; \h{larpix\_clk}'s
  \h{C\_RW\_REG3\_DEFAULT} is \h{0x00000000}, the value always built (the source's \h{0X00000004} was parsed as 0).
  The hijinks firmware starts here.
\item \textbf{Build option rot8} (7062e17): \h{pacman-rot8.xdc} copied onto \h{pacman.xdc} (tiles 1--8 on tiles
  2--10), as the 2024 rot8 builds were made.
\end{itemize}

\paragraph{Build.}  With the Vivado 2020.2 settings sourced: \h{tcl/recreate\_xpr.tcl}, then
\h{tcl/build\_xpr.tcl}.

\section{Excursions}

Each excursion is recorded by one \h{git merge -s ours} on 2425e67: the excursion's tip becomes the merge's second
parent, so all its commits keep their original hashes, while the tree never changes (nothing is combined).  The
first-parent log is the index; a single-branch clone holds every excursion.

\begin{table}[h]
\centering\small
\begin{tabularx}{\textwidth}{@{}lllllL@{}}
\toprule
Branch & Excursion & Owner & Base & Tip & What \\
\midrule
firmware & pacman-v1.2b-1MHz & Peter Madigan & v1.2b-0.0 & ba2ef0a & 1 MHz MCLK for the 1v2b board \\
firmware & pacman-v1.3-1MHz & Peter Madigan & v1.3-0.2 & 38af5a7 & fixed 1 MHz MCLK \\
firmware & pacman-v1.3-5MHz & Peter Madigan & v1.3-0.2 & 8a7e990 & fixed 5 MHz MCLK \\
software & pacman-v1.2b (v1.2b-0.1) & Peter Madigan & v1.2b-0.0 & 6d70ed2 & maintenance release (packaging) \\
software & pacman-1v4-petalinux20.2 & Armin Karcher & v1.4-0.2 & 5e67360 & PetaLinux 2020.2 port, eMMC boot \\
software & legacy-pacman-1v4-pe & Michael Mulhearn & (the port) & 469c948 & + I2C fix for the MUX: the best 1v4
server \\
\bottomrule
\end{tabularx}
\caption{The excursions.  The fixed-MCLK lines are superseded by v1.3-1.0's adjustable MCLK.}
\end{table}

\section{Bitstreams and their reproduction}

\begin{longtable}{@{}>{\raggedright\arraybackslash}p{0.19\textwidth}>{\raggedright\arraybackslash}p{0.11\textwidth}>{\raggedright\arraybackslash}p{0.17\textwidth}>{\raggedright\arraybackslash}p{0.08\textwidth}>{\raggedright\arraybackslash}p{0.35\textwidth}@{}}
\toprule
Bitstream & Body & Commit & Vivado & Reproduction \\
\midrule
\endhead
v1.4-0.2 release & 1362d6b3 & 0681491 (tag) & 2018.3 & bitwise, the tag as committed (2026-10-13) \\
893 (1v3 boards) & 3194812d & a5f8e91 + build step & 2021.1 & bitwise from a fresh clone with the in-repo build
step (2026-10-31); reference read from a running 1v3 board \\
883 & 0ff4bf66 & 9fbc1dc & 2021.1 & bitwise from a fresh clone with the in-repo build step (2026-10-31) \\
894 (resethack 1.8) & ec9ecc71 & c98f6b0 + build step & 2021.1 & bitwise from a fresh clone with the in-repo
build step (2026-10-31); reference: the ``84''/r8hack boot images \\
SYNC FIX & 4fdfb7d2 & 3155d4a & 2020.2 & bitwise from a fresh clone with the in-repo build step (2026-11-01);
reference: LBNL 1v5\_20240731A\_ucd\_syncfix \\
SYNC FIX rot8 & 5fdf3d45 & 7062e17 & 2020.2 & bitwise (2026-10-14); reference: LBNL 1v5\_20240810A\_ucd\_rot8 \\
``74'' (2x2 image) & f0b74071 & --- & --- & not reproduced: built outside the repository, no project archive \\
\bottomrule
\end{longtable}

Pending: every bitstream above built once more from a fresh clone of the repository with its own scripts (in progress;
rot8, v1.4-0.2); the tags without an artifact (v1.3-0.0 to 0.3, v1.3-1.0, v1.4-0.0) built once and
their bodies recorded; v1.3-1.0 and v1.3-0.3 compared with 60209db9, the bitstream of a 2022-10-31 BOOT.BIN found
deleted on the 2x2 run-1 SD card, which would give the 1v3 line its release reference.

\section{Conventions}

\begin{itemize}
\item Historic commits keep their author, dates and hash.  The snapshots, build steps, build option and excursion
  merges are authored by augie; the hijinks 2.X restores by Michael Mulhearn.
\item Work brought in from elsewhere enters as a new commit, described by content: no cherry-pick lines, no
  repository-qualified branch names.  Bitstream hashes and bundle names are data and may appear.
\item A snapshot changes only the files its archive changes; a build step changes only build scripts and fixes
  without behavior change; each records its bitwise result in its message.
\end{itemize}

\section{Old branches}

Deleted on 2026-10-31, after a check that every commit they held is reachable from a kept branch or tag, or has its
tree there: \br{legacy-firmware-historic}, \br{working-on-legacy-firmware-buildable} (the earlier design with
build loops), the old excursions record, \br{hijinks-hw1v5-2020.2}, \br{legacy-pacman-1v4-pe},
\br{pacman-server-hijinks-2018.3}, and the tags \h{sync-fix} and \h{hijinks-2v8}.  The only commit let go
with a tree of its own is the first 2.8 restore (61e2fa4): the kept 2.8 plus the files its successor removes,
which equal 2425e67's.  Six further branches (the 2020.2 GPIO-only hwcheck, the 2024 v1.5 reorganization, the
2020.2 and 2023.2 hardware demos, a 3v4 driver fixup) are under discussion.

\section{Still open}

Freezing (the branches renamed without \br{working-on-}), and annotated tags on every state of interest, each
carrying its build result (toolchain, how it builds, configuration body and build-product hashes), under a prefix
of our own; a single \br{legacy} branch merging the line branches is under discussion.

\end{document}
